\subsection{The block decomposition for \texorpdfstring{$\zeta(2)$}{zeta(2)}}
\label{app:z2-blocks}

Conservation permits the four-block path of
Section~\ref{subsec:z23}. Its last block takes a finite
difference of order $11n$, leaving a polynomial part of degree at most
$13n-1$. The resulting binomial sums and their harmonic terms are given
in the supplementary notebook.

For $C_n=\prod_{u=9,10,11}\binom{un}{8n}$, put
\[
\begin{aligned}
 \rho(x,y)={}&\lfloor11x\rfloor-\lfloor y\rfloor-\lfloor11x-y\rfloor
             -3\lfloor8x\rfloor\\
 &+\sum_{u=9,10,11}
   \bigl(\lfloor ux+y\rfloor-\lfloor(u-8)x+y\rfloor\bigr),
 \qquad \phi(x)=\min_{0\le y<1}\rho(x,y).
\end{aligned}
\]
Legendre's formula and the divided-difference bounds give
\[
 \frac{C_nd_{12n}d_{13n}}{\mathcal D_n^{\mathrm{car}}}
 (p_n,q_n)\in\mathbb Z^2,\qquad
 \mathcal D_n^{\mathrm{car}}
 =\prod_{\sqrt{50n}<p\le13n}p^{\phi(n/p)}.
\]
Let $\mathcal D_n^{\mathrm{sh}}$ be the product of the primes
$p>\sqrt{50n}$ in the six open intervals
\[
 \left(\frac{13n}{2},8n\right),\quad
 \left(\frac{13n}{3},\frac{9n}{2}\right),\quad
 \left(\frac{13n}{4},\frac{10n}{3}\right),\quad
 \left(\frac{13n}{5},\frac{8n}{3}\right),\quad
 \left(\frac{13n}{6},\frac{9n}{4}\right),\quad
 \left(\frac{13n}{12},\frac{9n}{8}\right).
\]
For each such prime, the normalized sum reduces to
$\sum_{x\in\mathbb F_p}P(x)$ with $\deg P<p-1$, hence vanishes.
Thus, with $\mathcal D_n=\mathcal D_n^{\mathrm{car}}
\mathcal D_n^{\mathrm{sh}}$, the same integrality statement holds
after division by $\mathcal D_n$.

The prime number theorem gives
\[
 \lim_{n\to\infty}\frac1n
 \log\frac{d_{12n}d_{13n}}{\mathcal D_n}
 =25-\gamma-\frac{233}{120}<10.929906,\qquad
 \gamma=\int_{1/13}^{\infty}\frac{\phi(x)}{x^2}\,dx.
\]
The one-periodic profile $\phi$ has $94$ cells on $[0,1)$;
exact summation over $2000$ periods with a bounded tail gives
$\gamma>12.1284279747910$.
The shell calculations and integral bounds are reproduced in
the supraliminal material
combines them with the spectral bounds to obtain Lemma~\ref{lem:z2-growth}.
The supplied formalization and full written proof are at
\url{https://github.com/RamanujanMachine/zeta2_four_trajectory_matrices}.

\subsection{The arithmetic of the pure
\texorpdfstring{$\zeta(3)$}{zeta(3)} approximations}
\label{app:zeta3-measure}

Write
\[
 F_n(t)=\frac{\prod_{u=0}^{2}(t-(10+u)n)_{10n}}
 {(t+n)_{16n+1}(t+2n)_{14n+1}}
 =\sum_{j=n}^{17n}\left(\frac{a_j}{(t+j)^2}+\frac{b_j}{t+j}\right),
\]
and set
\[
 Q_n=\sum_ja_j,\qquad
 P_{3,n}=\sum_j\left(a_jH^{(3)}_{j+2n}
                       +\frac12b_jH^{(2)}_{j+2n}\right),\qquad
 C'_n=\frac{(15n)!^2}{(10n)!^3}.
\]
Gamma reflection and comparison of Laurent coefficients give
$C'_na_{2n+k}=(-1)^nB_{n,k}$ and
$C'_nb_{2n+k}=-(-1)^nB'_{n,k}$, including the simple-pole tails.
Hence $C'_n(Q_n,P_{3,n})=(-1)^n(\mathbf q_n,\mathbf p_{3,n})$,
so $P_{3,n}/Q_n=p_{3,n}/q_n$ for every $n\ge0$.
The error estimate of Lemma~\ref{lem:z34-selected-error} therefore applies.

The coefficient bounds $C'_na_j\in\mathbb Z$ and
$d_{16n}C'_nb_j\in\mathbb Z$ give
$6C'_nd_{19n}^3(Q_n,P_{3,n})\in\mathbb Z^2$.
For $p>29$, $p^2>30n$ and $p\nmid n$, a stronger common divisor
$\mathcal D'_n$ combines the factorial carry count with fourteen
first-order shells and the two second-order windows
\[
 \{n/p\}\in
 \left(\frac1{10},\frac2{19}\right)
 \cup\left(\frac{11}{12},\frac{15}{16}\right).
\]
The carry and shell counts are evaluated on the $270$ cells of the
Farey partition of order $29$, with periodic lifts.
The first-order terms cancel over $\mathbb F_p$; a trace identity
gives the additional factor in the two windows.
The resulting divisor satisfies, for $n\ge1$,
\[
 \frac{6C'_nd_{19n}^3}{\mathcal D'_n}(Q_n,P_{3,n})\in\mathbb Z^2,
 \qquad
 \liminf_{n\to\infty}\frac{\log\mathcal D'_n}{n}
 >30.458547243805.
\]
Together with
$\limsup n^{-1}\log|C'_nQ_n|<65.451536565866$, this gives
\[
 h<65.451536565866+57-30.458547243805<91.992990,
\]
as used in Theorem~\ref{thm:z23-measures}.

The sum identity and spectral calculations are in
the supplementary material the carry profiles, shell certificates, and density bounds are in
the supplementary material.

